\documentclass[11pt]{article}
\usepackage[margin=1in]{geometry}
\usepackage{amsmath,amssymb}
\usepackage{graphicx}
\usepackage{booktabs}
\usepackage{xcolor}

\setlength{\parskip}{6pt}
\setlength{\parindent}{0pt}

\title{Modeling Horse Race Outcomes --- HKJC}
\author{Project Methodology \& Production Notes}
\date{\today}

\begin{document}
\maketitle

\tableofcontents
\vspace{1em}

% =============================================================================
\section{Introduction}

This document covers the full methodology of the HKJC horse-racing probability
model, from first principles through the current production configuration. It
unifies three strands of work:

\begin{enumerate}\item \textbf{Modeling frameworks} --- binary logistic, exploded logit,
      softmax, and the market-anchored \emph{Offset} model that is the
      current production choice.
\item \textbf{Feature engineering} --- the 25 features used in production,
      why each is constructed the way it is, and what alternatives were
      tested and rejected.
\item \textbf{Betting \& results} --- Kelly bet sizing, bagging-based
      variance reduction, and cumulative performance over 2015--2026
      (including out-of-sample 2025 and partial 2026).
\end{enumerate}

The production recipe, in one line: \textbf{5-seed bagged Offset model
with 25 features (including the \texttt{body\_wt\_mkt\_residual} feature),
trained on all history prior to the test year, bet with Kelly
fraction 0.02 on a \$10M bankroll}.

Cumulative ROI over 2015--2026 at Kelly 0.02: \textbf{+332.5\% without
rebate, +880.1\% with 10\% rebate on race-losses $\geq$ \$10k},
with 9/12 and 11/12 positive years respectively.

% =============================================================================
\section{Horse Racing Context and HKJC Data}

\subsection{Race Structure}
Each race is classified by:
\begin{itemize}\item \textbf{Distance} --- 1000m, 1200m, 1400m, 1600m, 1650m, 1800m,
      2000m, 2200m. Distances vary by racecourse and day.
\item \textbf{Surface} --- \emph{Turf} or \emph{AWT} (all-weather track).
\item \textbf{Class tier} --- Class 1 (toughest) through Class 5, or Group races.
      Class is handicap-rating-banded.
\item \textbf{Racecourse} --- Sha Tin (ST, larger) or Happy Valley (HV,
      tight 1.2km circumference).
\end{itemize}

\subsection{Variables used}
For each horse in each race, HKJC publishes (pre- and post-race):
\begin{itemize}\item Identity: \texttt{horse\_id}, trainer, jockey.
\item Structural: \texttt{Draw} (stall position), \texttt{Act\_Wt} (weight carried),
      \texttt{Declar.\ Horse\ Wt.} (declared body weight), \texttt{Gear}.
\item Market: \texttt{Win Odds} (final). Implied probability
      $p_{\text{implied}} = 1 / \text{odds}$.
\item Outcome: \texttt{Pla.} (finish position), \texttt{Finish Time},
      \texttt{Running Position} (positions at 3--6 intra-race sections).
\end{itemize}

Additional data we scrape separately:
\begin{itemize}\item \textbf{Barrier trials} (pre-race gallops, 53k+ records).
\item \textbf{Sectional times} (per-section times and positions, 648k rows).
\item \textbf{Dividend information} for refunds and payout verification.
\end{itemize}

\subsection{Data hygiene}
\begin{itemize}\item Refunded runners (\texttt{WV}, \texttt{WV-A}, \texttt{WX}, \texttt{WX-A}, \texttt{WXNR})
      are dropped entirely.
\item Void finishers (\texttt{PU}, \texttt{DNF}, \texttt{DISQ}, \texttt{UR}, \texttt{FE}, \texttt{TNP})
      are coded 99 and excluded from training.
\item Races with gappy or duplicated positions are removed pre-training
      (see \texttt{remove\_tied\_entries}).
\end{itemize}

% =============================================================================
\section{Data Representation and Notation}

Each observation is indexed by $(\text{race\_id}, \text{horse\_id})$. Let
$R$ = number of races, $m_r$ = field size of race $r$, and
$N = \sum_r m_r$ = total observations.

Feature matrix $X \in \mathbb{R}^{N\times p}$; for observation $i$,
feature vector $x_i \in \mathbb{R}^p$ ($p = 25$ in production).

Outcomes:
\begin{itemize}\item Binary win indicator $y_i^{\text{win}} \in \{0,1\}$; exactly one horse per
      race has $y^{\text{win}}=1$.
\item Ordinal finish position $y_i^{\text{place}} \in \{1,\dots,m_r\}$.
\end{itemize}

% =============================================================================
\section{Modeling Approaches}

We developed three gradient-boosting objectives and evaluated a fourth
architecture (Offset) that became production.

\subsection{Binary Logistic-Loss Boosted Trees}

Each horse treated as a separate binary example. Given $F(x_i) \in \mathbb{R}$,
$p_i = \sigma(F(x_i))$ where $\sigma(z) = 1/(1+e^{-z})$.

\[
\mathcal{L}_{\text{logistic}}(F) \;=\; \sum_{i=1}^{N}
  \Bigl[-y_i^{\text{win}}\log p_i \;-\; (1-y_i^{\text{win}})\log(1-p_i)\Bigr].
\]

Gradient / Hessian: $\partial\ell/\partial F = p_i - y_i$,
$\partial^2\ell/\partial F^2 = p_i(1-p_i)$.

Ensemble: $F_T(x) = \sum_{m=1}^{T} \nu\, h_m(x)$,
each $h_m$ a regression tree of depth $\leq D$ trained on pseudo-residuals.

Post-hoc calibration: \texttt{CalibratedClassifierCV} (isotonic) on an
inner cross-validated fold, then per-race softmax normalization so
probabilities sum to 1 within a race.

\subsection{Exploded Logit (Plackett--Luce)}

Models the full ranking of horses in a race. For race $r$ with
latent scores $f_i = F(x_i)$ and true ranking positions $y_i^{\text{place}}$,
the likelihood of the observed ordering is:

\[
\mathcal{L}_r = \prod_{k=1}^{K}
  \frac{\exp(f_{\pi_k})}{\sum_{j \in \text{remaining at stage }k} \exp(f_j)}
\]

where $\pi_k$ is the horse in position $k$. With truncation at top-$n$
(winner-only when $n=1$, top-3 when $n=3$, full ordering when $n=-1$).

Negative log-likelihood, gradient, and Hessian are coded into a
custom XGBoost/LightGBM objective (\texttt{exploded\_logit\_obj}).

\subsection{Softmax (Winner-Only)}

Equivalent to exploded logit with $n=1$: only the winner's relative
probability is modeled. Trained with a custom softmax objective
(\texttt{softmax\_obj}).

Post-hoc calibration for softmax uses binned temperature scaling
(per field-size bin, optionally per surface).

\subsection{Offset Model (Production)}

\textbf{This is the current production model.} Instead of learning
the log-odds from scratch, the Offset model anchors predictions to the
market and learns only the \emph{residual}.

Let $\pi_i$ denote the normalized market implied probability for
observation $i$. Define the market log-odds
\[
\eta_i^{\text{mkt}} = \operatorname{logit}(\pi_i) = \log\frac{\pi_i}{1-\pi_i}.
\]

We train a binary-logistic XGBoost with $\eta_i^{\text{mkt}}$ as a fixed
\texttt{base\_margin}:
\[
p_i = \sigma\bigl(\eta_i^{\text{mkt}} + F(x_i)\bigr),
\]
where $F(\cdot)$ is the XGBoost ensemble. When $F(x_i)=0$, the model
reproduces the market. XGBoost's rounds adjust only $F(x_i)$, the
\emph{correction} to market.

\paragraph{Why this helps.}
\begin{enumerate}\item The market is an excellent calibrated probability estimator. Asking
      the model to re-derive it from scratch wastes tree capacity.
\item \textbf{Importance is evenly distributed across edge features}
      (gain $\approx 4$--$6$) rather than concentrated on the market
      anchor. This reveals which features are actually contributing
      residual signal.
\item Post-hoc isotonic calibration \emph{hurts} Offset --- the market
      anchor is already calibrated; reshaping the output distorts it.
      See Section \ref{sec:rejected}.
\end{enumerate}

\paragraph{Offset hyperparameters (deep\_slow, tuned).}
\begin{itemize}\item \texttt{learning\_rate=0.01}, \texttt{max\_depth=5}
\item \texttt{min\_child\_weight=20}
\item \texttt{subsample=0.7}, \texttt{colsample\_bytree=0.7}
\item \texttt{gamma=0.2}, \texttt{reg\_alpha=2.0}, \texttt{reg\_lambda=10.0}
\item \texttt{rounds=400}
\end{itemize}

Depth 5 is deeper than typical softmax configs because the model only
learns corrections --- deeper trees capture feature interactions without
overfitting the market signal.

% =============================================================================
\section{Production Feature Set (25 Features)}

All features are \emph{time-respecting}: any expanding / rolling
statistic uses only past races (shift-by-1 or expanding-minus-current).

\subsection{Market (2)}
\begin{itemize}\item \texttt{Implied\_Prob} $= 1/\text{Win Odds}$, renormalized per race.
\item \texttt{log\_implied\_prob} $= \log(\text{Implied\_Prob})$.
\end{itemize}

\subsection{Race card (4)}
\begin{itemize}\item \texttt{Racecourse} --- ST=1, HV=0.
\item \texttt{Track\_Turf} --- Turf=1, AWT=0.
\item \texttt{Draw} --- stall number 1--14.
\item \texttt{Act\_Wt} --- weight carried, lbs.
\end{itemize}

\subsection{Class \& Weight (3)}
\begin{itemize}\item \texttt{class\_change} $= \text{prev\_class} - \text{current\_class}$.
      The most important single feature by XGBoost gain (rank 1--2 nearly
      every year). Captures difficulty-level change.
\item \texttt{wt\_z} --- $z$-score of \texttt{Act\_Wt} within the current field.
\item \texttt{setup\_weight\_z} --- $z$-score of \texttt{Act\_Wt} vs horse's own
      expanding history of carried weights.
\end{itemize}

\subsection{Form features (4)}
\begin{itemize}\item \texttt{recent\_form} $= \text{avg}(1/\text{place})$ over last 3 races.
      Concave weighting: 1st $\to$ 1.0, 2nd $\to$ 0.5, 3rd $\to$ 0.33, tapering
      for back-of-pack. Empirically outperforms linear normalization.
\item \texttt{form\_vs\_career} = recent\_form $-$ career\_form. Trajectory signal.
\item \texttt{career\_beat\_odds} --- expanding mean of
      $(y_i^{\text{win}} - \pi_i)$ per horse. A per-horse market-residual.
\item \texttt{prev\_win} --- binary: won last race.
\end{itemize}

\subsection{Trainer features (3)}
\begin{itemize}\item \texttt{trainer\_track\_spec} --- trainer's win rate at this specific
      racecourse minus trainer's overall win rate. Only applied with
      min-runs thresholds (50 overall, 20 per track).
\item \texttt{tt\_gap} --- trainer $\times$ track actual WR minus market-implied WR.
\item \texttt{tr\_gap} --- trainer overall actual WR minus market-implied WR.
\end{itemize}

\subsection{Draw \& field interactions (2)}
\begin{itemize}\item \texttt{draw\_outside\_ST} --- binary: \texttt{Draw} $\geq 10$ \textbf{and} \texttt{Racecourse}=ST.
      Outside-draw penalty at Sha Tin (except ST 1000m straight, where it
      reverses --- the tree learns the reversal from other features).
\item \texttt{fav\_field\_size} = \texttt{is\_fav} $\times$ \texttt{field\_size}.
      Captures the small-field favourite overpricing effect.
\end{itemize}

\subsection{Barrier trial features (2)}
Market-orthogonal signals from pre-race gallops.
\begin{itemize}\item \texttt{bt\_avg\_early} --- expanding average of
      $\text{early\_pos}/\text{field\_size}$ in past barrier trials. Running-style signature.
\item \texttt{bt\_last\_behind} --- time behind winner in the most recent
      barrier trial.
\end{itemize}

\subsection{Sectional features (4)}
Computed from per-horse-per-section times.
\begin{itemize}\item \texttt{trail\_sect\_closing} --- avg closing speed vs race average
      (last-section speed differential), trailed across past races.
\item \texttt{trail\_sect\_peak\_at} --- average position (0=early, 1=late)
      of horse's peak sectional speed. Running-style signature.
\item \texttt{trail\_sect\_avg} --- overall sectional-speed deviation vs
      race average.
\item \texttt{trail\_sect\_best\_gain} --- biggest intra-race position gain
      in one section, averaged across races. Explosive-move capability.
\end{itemize}

\subsection{Body weight residual (1)}
\label{sec:body_wt_residual}

The newest and most impactful feature added this session.

\paragraph{Construction.} For each row, define
\[
\text{bucket} = (B_{\Delta w}, B_\pi),
\]
where $B_{\Delta w}$ bins $\text{body\_wt\_change}$ into
$\{<\!-10, [-10,-5], [-5,+5], [+5,+10], >+10\}$ lbs, and $B_\pi$ bins
\texttt{Implied\_Prob} into $\{<0.04, [0.04,0.08], [0.08,0.15], [0.15,0.25], >0.25\}$.
That yields up to $5\times 5=25$ buckets.

Time-respecting expanding statistics per bucket:
\[
\text{past\_wins}_i = \sum_{j<i,\ \text{bucket}(j)=\text{bucket}(i)} y_j^{\text{win}},
\qquad
\text{past\_implied}_i = \sum_{j<i,\ \text{bucket}(j)=\text{bucket}(i)} \pi_j.
\]

The feature is a shrunk residual:
\[
\texttt{body\_wt\_mkt\_residual}_i \;=\;
\frac{\text{past\_wins}_i - \text{past\_implied}_i}{\text{past\_implied}_i + k},
\qquad k=50.
\]

This is the \textbf{market-failure residual}: it measures, bucket by bucket,
the historical deviation between actual wins and market-implied wins.
Positive $\Rightarrow$ market systematically underprices this
(body\_wt\_change $\times$ implied\_prob) combination; negative $\Rightarrow$
systematically overprices.

\paragraph{Empirical signal.} The strongest effects are:
\begin{itemize}\item Favourites ($\pi > 0.15$) with moderate weight drops or gains:
      +2 to +3\% edge.
\item Horses after long rest ($>$60 days): weight \emph{gain} is positive
      (horse stronger).
\item Volatile-weight horses with big drops: +0.7\% edge.
\item Previous winners with weight drops: +1.3\% edge.
\end{itemize}

\paragraph{Why this specifically succeeded.} Earlier sessions tried many
``target market-failure residual'' features (e.g., draw-bias residual,
class-change residual); most hurt. Two distinguishing factors for body
weight:
\begin{enumerate}\item \texttt{Declar.\ Horse\ Wt.} is nearly orthogonal to market
      odds (correlation $\approx 0.01$). The market has minimal
      pricing of body-weight dynamics.
\item Effects are strongly non-monotonic and interaction-dependent
      (different direction by implied-prob bucket). A single raw
      feature cannot capture this; a bucketed residual can.
\end{enumerate}

Across 8 seeds: mean paired $\Delta = +34$pp cumulative ROI over the
baseline, 7/8 seeds positive, minimum-case ROI compressed from
$-33\%$ to $-3\%$.

% =============================================================================
\section{Architecture: 5-Seed Bagging}

\subsection{Method}

For each test year, train 5 independent Offset models with different
XGBoost random seeds, $s \in \{7, 17, 42, 123, 256\}$. Each model sees the
same data but makes different random choices in column subsampling,
row subsampling, and split tie-breaking. Predictions are averaged:
\[
\hat p_i^{\text{bagged}} = \frac{1}{5} \sum_{s=1}^{5} \hat p_i^{(s)},
\]
then renormalized per race so probabilities sum to 1.

\subsection{Effect}

\begin{itemize}\item \textbf{Log-loss unchanged} (mean $\approx 0.2452$ across seeds and
      bagged). Bagging does not improve prediction accuracy.
\item \textbf{Prediction variance reduced} --- average per-horse
      probability becomes more stable.
\item \textbf{Kelly amplification of noise} drops --- bagged bet sizes
      are more consistent race-to-race.
\item Cumulative ROI in a 3-meta-seed test: mean $+71\%$ vs single-seed
      mean $+49\%$ (over the 2015--2024 window). Min cum ROI across meta-seeds
      went from $-3\%$ (single-seed) to $+41\%$ (bagged). \textbf{Downside
      effectively eliminated.}
\item Historically losing years 2018 and 2019 flipped from negative to
      positive under bagging.
\end{itemize}

5 seeds is a good point on the diminishing-returns curve: further seeds
add marginal stability at $5\times$ compute per additional set.

% =============================================================================
\section{Kelly Bet Sizing}

\subsection{Multinomial Kelly}

For a race with field size $m$, probabilities $p_1,\dots,p_m$
(model), payoffs $b_1,\dots,b_m$ (decimal odds), Kelly maximises
expected log-wealth:
\[
\max_{w_1,\dots,w_m \geq 0,\ \sum w_k \leq 1}\ \sum_k p_k \log\!\left(1 - \sum_j w_j + w_k b_k\right).
\]

In practice we scale the optimal fractions by $\alpha$
(\emph{fractional Kelly}) to reduce variance:
$w_k \to \alpha w_k$. Production uses $\alpha = 0.02$.

Constraints in the \texttt{multinomial\_kelly} implementation:
\begin{itemize}\item \texttt{max\_per\_horse} $= 15\%$ of bankroll.
\item \texttt{max\_total} $= 30\%$ of bankroll per race.
\item Minimum bet \$10; sub-minimum stakes are rounded up or dropped.
\end{itemize}

\subsection{Rebate}

HKJC offers a rebate of 10\% on race-level losses at or above \$10k for
high-volume players. The simulation reports \emph{both} trajectories:
``without rebate'' (raw Kelly P\&L) and ``with rebate'' (adds
$0.1 \times \max(0,\text{race\_loss} - 10{,}000)$ back to the bankroll
at the end of each race where the condition triggers).

\subsection{Kelly sweep}

At Kelly fractions $\{0.01, 0.02, 0.03, 0.05\}$:

\begin{center}
\begin{tabular}{lrrrr}
\toprule
Metric & K=0.01 & K=0.02 & K=0.03 & K=0.05 \\
\midrule
Cum ROI (no rebate)   & +139.5\% & \textbf{+332.5\%} & +522.7\% & +644.8\% \\
Cum ROI (with rebate) & +260.3\% & +880.1\%          & +2027.3\% & +5701.8\% \\
Positive years (no rebate)   & 9/12 & \textbf{9/12} & 7/12 & 7/12 \\
Positive years (with rebate) & 11/12 & \textbf{11/12} & 10/12 & 9/12 \\
\bottomrule
\end{tabular}
\end{center}

Kelly 0.02 is the recommended production setting: 9/12 positive years
without rebate, 11/12 with rebate, and the best Sharpe-like balance.
Kelly 0.05 has higher mean but larger drawdowns in bad years.

% =============================================================================
\section{Feature-Engineering Methodology}
\label{sec:methodology}

This section captures the discipline that produced the current feature
set and the design principles learned along the way.

\subsection{Testing protocol}

\begin{enumerate}\item \textbf{Empirical bucket analysis first.} Before proposing a new
      feature encoding, measure the raw effect across buckets on
      historical data. Look for:
      \begin{itemize}
        \item Effect size ($|\text{edge}| \geq 0.5\%$ is a minimum bar).
        \item Year-by-year consistency (11+/14 years sameside is strong).
        \item Non-monotonicity or bucket-specific direction.
      \end{itemize}
\item \textbf{Paired-delta test across $\geq 3$ seeds.} Single-seed ROI
      differences are noise-dominated at the scale of single-feature
      changes. Require:
      \begin{itemize}
        \item Mean paired $\Delta$ $>$ 0 with all seeds same-sign, OR
        \item If mean is positive but some seeds negative, suspect
              regression-to-the-mean on lucky seeds.
      \end{itemize}
\item \textbf{Log-loss as primary predictive metric.} LL is stable across
      seeds (std $\sim 5\times 10^{-5}$); ROI is high-variance (std
      $\sim 15$--$35\%$). If LL is unchanged, the feature doesn't
      improve predictions --- ROI swing is Kelly amplification of noise.
\item \textbf{Multi-year ROI as production-relevant metric.} LL is what
      we're modeling, ROI is what we care about. Both must move to
      declare a win.
\end{enumerate}

\subsection{Design principles}

\begin{enumerate}\item \textbf{Target market-failure residuals, not market-known structure.}
      The market already prices known effects (draw bias, class drops,
      favourite-longshot). Encoding these as features duplicates
      information already in the offset's \texttt{base\_margin} and adds
      only noise.

      \emph{Example failure:} \texttt{draw\_bias\_loc} was a time-respecting
      bucketed lift feature for draw. Empirically accurate; in-model it
      cost $-21$pp cumulative ROI because the market already knows.
\item \textbf{Prune low-importance leaf features; preserve backbone
      categoricals.} Low gain-importance can mean two things:
      \begin{itemize}        \item Genuine dead weight (leaf feature with no interaction
              role) --- safe to prune.
        \item Backbone categorical (\texttt{Racecourse}, \texttt{Track\_Turf})
              that enables \emph{other} features' interactions ---
              dropping crashes the model.
      \end{itemize}
      Rule: if any other feature's computation or interaction depends
      on this one, do not prune.
\item \textbf{Don't re-encode top-importance features.} If a feature
      consistently ranks 1--3 in gain, XGBoost has already found
      optimal splits. Adding a derived version (e.g.,
      \texttt{class\_change\_residual} built from \texttt{class\_change})
      wastes tree capacity on correlated information.
\item \textbf{Shape matters for form encodings.} \texttt{recent\_form}
      using $1/\text{place}$ (concave, emphasising wins) outperforms
      linear normalization by $(\text{field\_size} - \text{place} + 1)/\text{field\_size}$.
      Horse quality is better signalled by top finishes than by
      average position.
\item \textbf{Re-engineering payoff is concentrated at mid-to-low
      importance.} Top features are already optimal; dead features
      are dead. The sweet spot for investment is features at gain-rank
      5--15 that have known market-failure structure.
\end{enumerate}

\subsection{What was tried and rejected}
\label{sec:rejected}

Full list of re-engineering attempts that failed, saved here so future
sessions don't repeat the work:

\begin{enumerate}\item \texttt{draw\_bias\_loc}: bucketed draw-bias residual. $-21$pp.
      \emph{Reason:} market already prices draw bias.
\item \texttt{class\_change\_residual}: bucketed residual on
      (current\_class $\times$ class\_change). $-38$pp across seeds.
      \emph{Reason:} \texttt{class\_change} is already rank 1 in gain;
      re-encoding adds overfitting surface without info.
\item \texttt{prev\_mkt\_residual} swap for \texttt{prev\_win}: neutral
      on LL, high-variance ROI. Mean $+9$pp but one seed crashed.
      \emph{Kept \texttt{prev\_win} binary.}
\item Form normalization $(\text{field\_size} - \text{place} + 1)/\text{field\_size}$
      replacing $1/\text{place}$: $-32$pp.
      \emph{Reason:} concave weighting of top finishes is the correct
      shape.
\item LOO on odds-dependent features (\texttt{fav\_field\_size},
      \texttt{career\_beat\_odds}, \texttt{tt\_gap}, \texttt{tr\_gap}):
      every removal hurt $-20$ to $-40$pp.
\item Sectional consolidation (4 $\to$ 2 or 1): $-22$ to $-47$pp.
      \emph{Reason:} all 4 features carry complementary information.
\item Pruning \texttt{Racecourse} or \texttt{trail\_sect\_closing}: $-30$
      to $-66$pp.
      \emph{Reason:} backbone for downstream interactions.
\item Isotonic post-hoc calibration on Offset: $-141$pp.
      \emph{Reason:} market anchor via \texttt{base\_margin} is already
      calibrated by construction; isotonic distorts the output.
\item Softmax + isotonic calibration (with \texttt{log\_implied\_prob}):
      $-91$pp.
      \emph{Reason:} Softmax must learn the market anchor from scratch;
      fundamentally weaker than hard-anchoring via base\_margin.
\item Bagged Offset + Softmax ensemble (50/50): $-62$pp.
      \emph{Reason:} blending a strong model with a weaker one yields a
      mediocre model.
\item Kelly $\alpha = 0.10$: catastrophic loss ($-92\%$ cumulative).
      Volatility drag compounds across bad years.
\item Adding raw body weight features (\texttt{body\_wt\_z} +
      \texttt{body\_wt\_change}): $-54$pp. \emph{Only the bucketed
      residual form works.}
\item \texttt{body\_wt\_after\_win} targeted binary: $-48$pp. Too sparse.
\end{enumerate}

\subsection{What worked}

\begin{enumerate}\item \textbf{Prune2} --- dropped \texttt{top3\_count\_5} and
      \texttt{field\_form\_cv} (lowest-gain leaf features).
      $+69$pp paired $\Delta$, 3/3 seeds positive.
\item \textbf{\texttt{body\_wt\_mkt\_residual}} --- bucketed body-weight
      market-failure residual (Section \ref{sec:body_wt_residual}).
      $+34$pp paired $\Delta$ across 8 seeds.
\item \textbf{5-seed bagging} --- averaging Offset predictions across
      independent seeds. $+22$pp mean, downside compressed.
\item \textbf{Higher Kelly safe with bagging} --- sweet spot shifted
      from $0.01$--$0.02$ (single-seed) to $0.02$--$0.03$ (bagged).
\end{enumerate}

% =============================================================================
\section{Implementation Details}

\subsection{Preparing group pointers}

XGBoost requires group sizes (one integer per race, giving the number
of horses). Our helper \texttt{prepare\_group\_sizes(X)} produces this
from a \texttt{MultiIndex}-ed DataFrame.

\subsection{Handling ties and missing positions}

For exploded-logit training, ties are broken by averaging the
negative-gradient penalty across tied horses. Missing positions (void
finishers) are excluded before any shift/expanding computation.

\subsection{Time-respecting cross-validation}

Production uses expanding-window temporal split: train years = all
years strictly before the test year. No fold contains future data.
Separate calibration year (the last train year) is carved out only
when \texttt{--calibrate} is explicitly set; not used for Offset.

% =============================================================================
\section{Results: Year-by-Year (2015--2026)}

\subsection{Market vs model log-loss}

Model beats market log-loss 10 of 10 years in the original test
(2015--2024), mean $\Delta$ = $-0.00089$ (excluding 2024 outlier: $-0.00047$).
The absolute improvement is small ($\sim 0.2\%$ relative) but consistent.

\subsection{Year-by-year ROI, Kelly 0.02}

Bagged Offset, per-year reset to \$10M, across 12 test years:

\begin{center}
\begin{tabular}{lrr}
\toprule
Year & ROI no-rebate & ROI with-rebate \\
\midrule
2015 & +82.35\% & +106.97\% \\
2016 & +68.03\% & +85.44\% \\
2017 & +2.26\% & +11.46\% \\
2018 & +0.21\% & +8.53\% \\
2019 & +7.81\% & +15.65\% \\
2020 & $-12.72\%$ & $-5.75\%$ \\
2021 & +24.78\% & +33.27\% \\
2022 & +12.72\% & +18.54\% \\
2023 & +7.58\% & +12.75\% \\
2024 & $-4.98\%$ & +1.49\% \\
2025 & $-2.83\%$ & +0.94\% \\
2026$^*$ & +4.77\% & +6.11\% \\
\midrule
Cumulative & +332.5\% & +880.1\% \\
Positive years & 9/12 & 11/12 \\
\bottomrule
\end{tabular}
\end{center}
$^*$2026 is partial (approximately 2{,}500 races vs 9{,}000 full year).

\subsection{Observations}

\begin{itemize}\item \textbf{2015--2016} produced the bulk of the cumulative alpha.
      Market LL in those years was low (market was sharp) --- the
      model still beat it handily.
\item \textbf{2020} is universally the worst year. Plausibly
      COVID-related disruptions (empty stadiums, jockey pool changes,
      betting-volume shift to online).
\item \textbf{Recent years (2023--2025)} show smaller edges, consistent
      with the ``market-efficiency edge decay'' hypothesis: as more
      sophisticated bettors arbitrage the same signals, our
      market-failure residuals shrink.
\item \textbf{Out-of-sample 2025 and 2026}: slight no-rebate loss in
      2025, positive in 2026 partial. Both positive with rebate. The
      model appears to generalize to post-training periods.
\end{itemize}

% =============================================================================
\section{Practical Concerns}

\subsection{Edge decay over time}

The market-failure features (\texttt{body\_wt\_mkt\_residual},
\texttt{tt\_gap}, \texttt{tr\_gap}, \texttt{career\_beat\_odds}) are
\emph{structural} residuals that automatically shrink as the market
prices their content. The feature still computes correctly year-by-year
(because it uses expanding history), but the magnitude of
$\text{past\_wins} - \text{past\_implied}$ per bucket decreases, so
Kelly bets smaller, so ROI shrinks.

This is the structural reason 2015--2016 were outsized: the market had
not yet priced certain patterns.

\subsubsection*{Truncated-training experiment (tested, rejected)}

To test whether the drift was severe enough to justify reducing the
training window, we backtested three windows across 2015--2026 at
three seeds:

\begin{center}
\begin{tabular}{lrrr}
\toprule
Window & Mean cum no-rebate & With rebate & Seeds no-rebate \\
\midrule
Full (2010--$t{-}1$) & \textbf{+103.34\%} & \textbf{+382.42\%} & +26.6\%, $-$18.9\%, +302.4\% \\
Last 5 years & $-71.37\%$ & +18.26\% & $-$69.2\%, $-$72.7\%, $-$72.2\% \\
Last 3 years & $-85.06\%$ & $-$1.42\%  & $-$67.2\%, $-$91.1\%, $-$96.9\% \\
\bottomrule
\end{tabular}
\end{center}

\textbf{Truncation catastrophically hurt overall.} Full history wins by
a huge margin on cumulative ROI.

However, the per-year breakdown shows \emph{recency-specific} gains
hidden inside the global failure:
\begin{itemize}\item 2024: last\_3 was $+7.92\%$ vs full $-4.46\%$ (flipped positive).
\item 2025: last\_3 was $+21.00\%$ vs full $-7.03\%$ (large flip).
\item 2026: last\_3 was $+3.12\%$ vs full $-0.30\%$ (flipped positive).
\end{itemize}

So the edge-decay story has \textbf{partial support}: very recent
years DO benefit from a shorter training window. But the cumulative
cost in 2015--2021 (where full history is massively better) dwarfs the
recency gain.

\subsubsection*{Implications for future work}

\begin{itemize}\item Simple truncation is too aggressive. Too much signal lives in
      older data.
\item \textbf{Recency-weighted training} --- weight recent races more
      ($w = \exp(-\lambda \cdot \text{years\_ago})$) without discarding
      older data --- is the natural next attempt. Untested.
\item \textbf{Full-plus-truncated ensemble} --- average predictions from
      a full-history model and a last-3-year model. Untested.
\item Regular re-screening of feature importance by year to catch
      decaying features early.
\end{itemize}

\subsection{Rebate dependency}

Before body\_wt + bagging: with-rebate was $+78\%$ vs no-rebate $+14\%$
--- almost all the alpha came from rebate. After body\_wt + bagging:
no-rebate cumulative is $+332.5\%$, with-rebate $+880.1\%$ --- alpha
is genuine, rebate adds significant further compounding but isn't
load-bearing.

Critical for production: the strategy works without rebate, so is
deployable for any player, though the tail compounds much higher with
rebate.

\subsection{Bet-time feature stability}

Several features (\texttt{fav\_field\_size}, \texttt{career\_beat\_odds},
\texttt{tt\_gap}, \texttt{tr\_gap}, \texttt{body\_wt\_mkt\_residual})
depend on odds observed pre-race. At actual bet time, odds can shift
due to late market flow. LOO tests showed every odds-dependent feature
earns its keep, so we keep them despite the fragility risk; production
deployment should re-run inference shortly before bet placement.

\subsection{Strategic jockey behaviour}

Several features assume independent race-to-race horse behavior
(\texttt{career\_beat\_odds}, form aggregates). If connections
strategically rest a horse's form for later, these backward-looking
features miss the signal. We have no observable signal for this
intentionality beyond \texttt{Gear} and \texttt{class\_change}.

% =============================================================================
\section{Reproducing the Pipeline}

\paragraph{Step 1: Build the feature parquet.}
\begin{verbatim}
python data-processing/prep-data.py
\end{verbatim}
Loads raw CSVs and JSONs, merges horse history, computes all 25
production features (including the body-weight market-failure residual),
and writes \texttt{data/processed/race\_features.parquet}.

\paragraph{Step 2: Run the production backtest + viz.}
\begin{verbatim}
python model/run_production.py
\end{verbatim}
Trains 5-seed bagged Offset for each test year 2015--2026, simulates
Kelly ${0.01, 0.02, 0.03, 0.05}$ bankroll trajectories, and writes
\texttt{model/results/bankroll\_viz.html} --- an interactive
visualization with per-Kelly and per-year tabs.

\paragraph{Streamlit UI.}
\begin{verbatim}
cd model && streamlit run app.py
\end{verbatim}
Defaults on the Training page are Offset + 5-seed bagging + the tuned
deep\_slow hyperparameters. The Bankroll page defaults to Kelly 0.02
and \$10M starting bankroll.

% =============================================================================
\section*{Appendix: Notation reference}

\begin{tabular}{ll}
$\pi_i$ & market implied probability for observation $i$ \\
$\sigma(\cdot)$ & sigmoid function \\
$F(x)$ & XGBoost ensemble output (log-odds correction in Offset) \\
$\eta_i^{\text{mkt}}$ & $= \operatorname{logit}(\pi_i)$, market log-odds \\
$y_i^{\text{win}}$ & binary win indicator \\
$y_i^{\text{place}}$ & finish position \\
$m_r$ & field size of race $r$ \\
$N = \sum_r m_r$ & total observations \\
$k$ & shrinkage prior for market-failure residuals (= 50) \\
$\alpha$ & Kelly fraction \\
\end{tabular}

\end{document}
